\subsection{A false lifting assertion in a modern auxiliary source}

\begin{proposition}[D-031: \status{Refuted}; the unrestricted lifting claim of Caradot--Lin]
\label{sa08:D031}
Proposition~3.10 of Caradot--Lin,
\emph{Differential graded algebras with divided powers and
homotopy Lie algebras}, arXiv:2511.22614v4,
pp.~19--20, is false with its stated hypotheses.
This concerns that proposition and does not refute
the distinct resolution comparison of their Theorem~3.15.
\end{proposition}

\begin{proof}
The source defines a free complex in \S2.1, p.~6,
by requiring that \(M^i\) is ``a free \(A\)-module for
all \(i\).'' Its Proposition~3.10 assumes a surjective
map on PD dg ideals and claims that a map from any
such termwise-free complex lifts as a chain map, hence
extends to a map of free PD dg algebras
\cite[\S2.1, p.~6; Proposition~3.10, pp.~19--20]{D_CaradotLin2026}.
Termwise freeness does not justify that chain-map lift.

Let \(k\ne0\) be any commutative ring. Define cochain
complexes
\[
 E=\bigl(ku\xrightarrow{\mathrm{id}}kv\bigr),
       \qquad |u|=-2,\quad |v|=-1,
 \qquad
 M=k\varepsilon,\quad |\varepsilon|=-2,\quad d\varepsilon=0.
 \tag{D-31a}
\]
Give \(k\oplus E\) and \(k\oplus M\) square-zero
multiplication on their displayed ideals. On every
even ideal element set \(\gamma_1(a)=a\) and
\(\gamma_n(a)=0\) for \(n\ge2\).
These are PD dg algebras: the PD sum and product
identities hold because products of ideal elements
are zero, and differential compatibility holds for
the same reason. In particular this example lies
entirely in nonpositive cochain degrees.

The map
\[
 \pi:k\oplus E\longrightarrow k\oplus M,
 \qquad \pi(u)=\varepsilon,\quad\pi(v)=0,
 \quad\pi|_k=\mathrm{id},
 \tag{D-31b}
\]
is a PD dg algebra morphism and is degreewise surjective
on the ideals. Take the termwise-free complex \(X=M\)
and the identity map \(\rho:X\to M\).
Any degreewise lift of \(\rho\) must send
\(\varepsilon\) to \(u\). But \(du=v\ne0\), while
\(d\varepsilon=0\). Hence no chain-map lift exists.

Equivalently, the free PD dg algebra on the even
generator \(\varepsilon\) has a PD dg map to
\(k\oplus M\); a lift to \(k\oplus E\) would have
to send its closed generator to \(u\), which is not
closed. Thus the conclusion of Proposition~3.10 itself,
not merely the displayed intermediate argument,
fails in this example.
\end{proof}

The exact finite certificate
\texttt{code/D\_lifting\_counterexample.py}
records the differential and projection matrices over
\(\Z\) and verifies \(\pi d=d\pi\).
The lift equations are
\[
 \begin{pmatrix}1\\1\end{pmatrix}c
       =\begin{pmatrix}1\\0\end{pmatrix}.
\]
Subtracting the second row from the first gives
\(0=1\), so the contradiction holds over every nonzero
coefficient ring. The script also exhausts the possible
scalars over \(\F_2\) and \(\F_3\). Its matrices and
outputs are saved in
\texttt{logs/D\_lifting\_counterexample.json}.

\paragraph{Valid lifting versions.}
There are two precise corrections. A projective object
\(X\) of the abelian category of complexes lifts through
every epimorphism of complexes, by the definition of
projectivity. Alternatively, suppose that \(X\) has
projective terms and is K-projective, and that the map
\(\pi:I'\to I\) of complexes of PD ideals is a
surjective quasi-isomorphism. Then every chain map
\(\rho:X\to I\) lifts to a chain map \(X\to I'\).

Indeed choose a degreewise lift \(g\), using the
projectivity of the terms of \(X\). Its defect
\(\delta=d_{I'}g-gd_X\) lands in \(K=\ker\pi\)
and is a degree-one cocycle of
\(\operatorname{Hom}^*(X,K)\). The kernel \(K\) is
acyclic. K-projectivity makes this Hom complex
acyclic, so \(\delta=d(h)\) for a degree-zero map
\(h:X\to K\). Then \(g-h\) is the desired chain
map lift. The free PD algebra property extends it
to a PD dg algebra map. A bounded-above cochain
complex of projective modules is K-projective, as
the usual descending construction of a nullhomotopy
against an acyclic complex proves. These corrected
versions retain the hypothesis needed to remove
the chain-map defect.

\paragraph{Scope of the rejected lead.}
The map (D-31b) is not a quasi-isomorphism: its source
has no negative cohomology, whereas its target has
\(k\) in degree \(-2\). Theorem~3.15 of the same
source has an additional target-resolution hypothesis
and constructs lifts of cycles through boundaries
\cite[Theorem~3.15, pp.~23--25]{D_CaradotLin2026}.
The counterexample above does not address that theorem.
It shows precisely why Proposition~3.10 cannot be
used to remove the simplicial compatibility requirement
in the integral C.3 problem. The present paper uses
the explicit bounded comparisons in (B-C1)--(B-C3)
and Statement~\ref{sa08:D030}, with their acyclicity
hypotheses, instead of this unrestricted lifting claim.
